\par\Needspace{6\baselineskip}\medskip
\hypertarget{native-124-article091}{}\pdfbookmark[2]{91.}{article-091}
\begin{center}91.\end{center}
In Gleichung 12 ist in dem hier betrachteten Fall offenbar das obere Vorzeichen zu wählen. Durch Einführung der Größe $\xi$ erhält diese Gleichung die Form
\[
m=(l+x)^{\frac12}+\dfrac{(l+x)^{\frac32}}{\frac34-\frac{9}{10}(x-\xi)}\mathpunct{\cdot}
\]
Setzen wir also $\BareRoot{(l+x)}=\dfrac{m}{y}\mathpunct{\raisebox{.45ex}{,}}$ und
\[
\mathrm{14)}\qquad \dfrac{mm}{\frac56+l+\xi}=h,
\]
so ergibt sich nach der Umformung
\[
\mathrm{15)}\qquad h=\dfrac{(y-1)yy}{y+\frac13}\mbox{\HistoricalDisplayNoteMark}.
\]
\noindent\HistoricalDisplayNoteText{Handschriftliche Bemerkung:]
\[
y-1=\dfrac{l+x}{\frac34-\frac{9}{10}(x-\xi)}
\]
\[
y=1+\frac{10}{9}h-\frac{110}{81}hh+\frac{2410}{729}h^3-\text{usw.}
\]
\[
yy=1+\frac{20}{9}h-\frac{40}{27}hh+\frac{2620}{729}h^3-\text{usw.}
\]
\[
\dfrac{1}{yy}=1-\frac{20}{9}h+\frac{520}{81}hh-\frac{5140}{243}h^3+\text{usw.}
\]
}Dürfen wir also $h$ als bekannte Größe betrachten, so lässt sich daraus $y$ durch die kubische Gleichung bestimmen, und anschließend gilt
\[
\mathrm{16)}\qquad x=\dfrac{mm}{yy}-l.
\]
Obwohl $h$ die noch unbekannte Größe $\xi$ enthält, darf man diese in der \SourceLeaf{125}ersten Näherung vernachlässigen und für $h$ den Wert $\dfrac{mm}{\frac56+l}\mathpunct{\raisebox{.45ex}{,}}$ ansetzen, denn $\xi$ ist in dem betrachteten Fall gewiss stets sehr klein. Aus den Gleichungen 15 und 16 erhält man dann $y$ und $x$. Aus Tafel III entnimmt man zu $x$ den Wert von $\xi$ und erhält damit nach Formel 14 den korrigierten Wert von $h$. Wiederholt man mit diesem Wert dieselbe Rechnung, so erhält man korrigierte Werte für $y$ und $x$. Diese weichen in der Regel so wenig von den vorigen ab, dass der erneut aus Tafel III entnommene Wert von $\xi$ mit dem ersten übereinstimmt. Andernfalls muss man die Rechnung wiederholen, bis sich dieser Wert nicht mehr ändert. Sobald $x$ bestimmt ist, erhält man $g$ aus der Formel $\sin\tfrac12 g^2=x$.

Diese Regeln gelten für den ersten Fall, in dem $\cos f$ positiv ist. Im anderen Fall, in dem er negativ ist, setzen wir $\BareRoot{(L-x)}=\dfrac{M}{Y}$ und
\[
\mathrm{14}^{\ast}\mathrm{)}\qquad \dfrac{MM}{L-\frac56-\xi}=H,
\]
womit Gleichung $12^{\ast}$ nach der Umformung in die folgende übergeht
\[
\mathrm{15}^{\ast}\mathrm{)}\qquad H=\dfrac{(Y+1)YY}{Y-\frac13}\mathpunct{\cdot}
\]
Mit dieser kubischen Gleichung lässt sich also $Y$ aus $H$ bestimmen. Daraus ergibt sich wiederum $x$ nach der Gleichung
\[
\mathrm{16}^{\ast}\mathrm{)}\qquad x=L-\dfrac{MM}{YY}\mathpunct{\cdot}
\]
In der ersten Näherung erhält $H$ den Wert $\dfrac{MM}{L-\frac56}\mathpunct{\raisebox{.45ex}{;}}$ mit dem Wert von $x$, den man daraus nach den Gleichungen $15^{\ast}$ und $16^{\ast}$ gewinnt, entnimmt man Tafel III den Wert von $\xi$. Damit erhält man nach Formel $14^{\ast}$ den korrigierten Wert von $H$, mit dem man die Rechnung in gleicher Weise wiederholt. Schließlich bestimmt man aus $x$ den Winkel $g$ ebenso wie im ersten Fall.
